\begin{figure}[h]
\centering
\begin{tikzpicture}[
  node distance=0.3cm and 0.12cm,
  every node/.style={\footnotesize},
  box/.style={
    draw=black!60, rounded corners=3pt, fill=#1,
    minimum height=0.5cm, text width=1.7cm, align=center,
    font=\footnotesize
  },
  box/.default={blue!6},
  arr/.style={-stealth, thick, black!55},
  lbl/.style={font=\footnotesize\bfseries, text=black!70},
]
  % Active row
  \node[lbl] (la) {Active};
  \node[box=orange!8, right=0.25cm of la]  (a1) {Blurred Image\\+ Question};
  \node[box, right=of a1]         (a2) {Model\\Answer};
  \node[box=violet!10, right=of a2] (a3) {Judge\\Admits?\\Fabricates?\\Correct?};

  \draw[arr] (a1) -- (a2);
  \draw[arr] (a2) -- (a3);

  % Passive row
  \node[lbl, below=0.55cm of la] (lp) {Passive};
  \node[box=orange!8, right=0.25cm of lp]  (p1) {Blurred Image\\+ ``Describe''};
  \node[box, right=of p1]         (p2) {Model\\Description};
  \node[box=violet!10, right=of p2] (p3) {Judge\\Mentioned?\\Admits?\\Fabricates?\\Correct?};

  \draw[arr] (p1) -- (p2);
  \draw[arr] (p2) -- (p3);
\end{tikzpicture}
\caption{Active and passive evaluation protocols for selective blur. Active probes ask targeted questions; passive probes use open-ended description prompts. Both measure admission and fabrication behaviour.}
\label{fig:uml-active-passive}
\end{figure}
